% Source: Independent derivation and implementation.
\subsection{二维前缀和}
\label{other:prefix-2d}
\AlgorithmInfo
    {静态矩阵中多次查询子矩形元素和。}
    {建表 $O(nm)$，查询 $O(1)$，空间 $O(nm)$。}
    {行列均从 1 开始，输入带第 0 行／列占位；所有元素绝对值之和 $\le4\times10^{18}$。}
\subsubsection*{原理}
令 $s_{i,j}=\sum_{x=1}^i\sum_{y=1}^j a_{x,y}$，第 0 行和第 0 列均为 0。按行列递增计算
\[
 s_{i,j}=a_{i,j}+s_{i-1,j}+s_{i,j-1}-s_{i-1,j-1}.
\]
上方与左方相加会重复左上区域，故减一次。闭矩形 $[x_1,x_2]\times[y_1,y_2]$ 的和为
\[
 s_{x_2,y_2}-s_{x_1-1,y_2}-s_{x_2,y_1-1}+s_{x_1-1,y_1-1}.
\]
\subsubsection*{使用教程}
\InlineCode{Prefix2D::build(a)} 建表并覆盖原状态，a 必须为规则矩形且至少含占位行／列。

\InlineCode{query(x1,y1,x2,y2)} 查询合法非空闭矩形，先行后列。更新 a 后应重新 build，不能直接继续使用旧表；查询中的中间加减提升到 128 位，结果仍是 long long。
\begin{lstlisting}[style=mtycpp,numbers=none]
vector<vector<long long>> a={{0,0,0,0},{0,1,2,3},{0,4,5,6}};
Prefix2D pre;pre.build(a);
long long sum=pre.query(1,2,2,3);
\end{lstlisting}
\needspace{7\baselineskip}
\subsubsection*{模板代码}
\CodeFile{code/others/prefix_2d.cpp}
